\section{\large\textbf{Conclusion}}
In this study, the influence of the directivity factor on binaural parameters and, consequently, on the ability to localize sound sources, has been determined. Although no significant differences were found overall, an improvement in localization ability was observed for the source with a directivity factor of 4, in the 400 and 500 Hz bands. This finding regarding the impact of source directivity suggests that this source modification significantly enhances localization in this range of the spectrum; however, the accuracy rate for other sources remained high (over 80\%), indicating that this improvement is not essential. Above 630 Hz, all sources exhibited similar behavior, with no one source standing out. For the 250 Hz band, localization ability decreased considerably, leading to the conclusion that localization difficulty increases at this frequency.

A considerable source of error may stem from the geometry of the directivity balloons used for each source. While the directivity 1 source had a perfectly omnidirectional directivity, no information on the spatial energy distribution was available for the other sources. This could result in concentrated lobes that increase the perceived energy at the measurement point, potentially producing unexpected results, such as the reduction in correct responses between directivities 1 and 2.

It would be beneficial to extend the study to even lower frequencies through simulations based on the wave equation. This approach could help determine the limits of localization ability and assess whether directivity influences lower frequencies. However, achieving directivity factors of 4 in the bands where localization improvements were noted remains impractical or prohibitively expensive. Given that this improvement is slight, pursuing this level of directivity would likely not be worthwhile.

Finally, conducting this study in situ would be valuable to assess whether differences from the current findings emerge. This could enable validation of the results and allow for testing within simulated living room geometries to investigate the impact on perception further.

It would be of interest to conduct a study on the localization of impulsive sounds in which low frequencies are reproduced through a fixed omnidirectional system (such as a subwoofer), while mid-high frequencies are reproduced through sources that are meant to be localized, using the crossover frequency as a parameter. This approach would complement the findings of the current study and allow for the consideration of such stimuli in the design of an electroacoustic system. According to previous studies, room geometry should be considered in these cases.
